The \wikilink{s29--b002}{Noetherian finite-generation criterion} can be
tested on prime ideals alone. All rings here are commutative and unital.

\begin{theorem}\label{cohen-criterion}
A ring $R$ is Noetherian if and only if every prime ideal is finitely generated.
\end{theorem}
\begin{proof}\label{cohen-criterion-proof}
Only the reverse implication needs proof. Suppose some ideal is not
finitely generated. The union of a nonempty chain of such ideals cannot
be finitely generated: its finitely many proposed generators would lie
in one chain member, making that member equal to the union and finitely
generated. Zorn's lemma gives an ideal $I$ maximal among the ideals that
are not finitely generated. It is proper since $R=(1)$.

Suppose $ab\in I$ with $a,b\notin I$. Both $I+(a)$ and
$(I:a)=\{r:ra\in I\}$ strictly contain $I$, the latter because it
contains $b$. By maximality both are finitely generated. Write
\[
 I+(a)=(a,u_1,\ldots,u_r),\quad u_i\in I,
 \qquad (I:a)=(v_1,\ldots,v_s).
\]
The $u_i$ may be obtained by subtracting multiples of $a$ from any
finite generating set. If $x\in I$, express $x=ca+\sum d_iu_i$.
Then $ca\in I$, so $c\in(I:a)$, proving
$I=(u_1,\ldots,u_r,av_1,\ldots,av_s)$.
All these generators lie in $I$, a contradiction. Therefore $I$ is
prime, contradicting the assumed finite generation of prime ideals.
\end{proof}
\begin{remark}\label{cohen-scope}
The maximal object is maximal among non-finitely-generated ideals, not
necessarily a maximal proper ideal. Neither a domain hypothesis nor
prior Noetherianity is used in the reverse implication.
\end{remark}
